% ch05: 原文 GLDraft_D.tex 行 1034–1151；§3
% 只译散文（标题、正文、图题、脚注）；公式、\label、\ref、\eqref、\cite、图文件名保持不变。

  \section{Hitchin 哈密顿量的本征值}\label{hitchval} 

\subsection{中心平凡的情形}\label{centertriv}

现在我们可以开始推出电–磁对偶的一些有趣推论。

一旦像我们在第 \ref{dualco} 节所做的那样，把 $\B_\cc$ 的 $B$-模型对偶认定为 $\B_\op$，就立刻得到 $\A=\Hom(\B_\cc,\B_\cc)$ 的对偶描述：它是 $\Hom(\B_\op,\B_\op)$。由于 $\B_\op$ 是支撑在 $L_\op$ 上的秩 1 拉格朗日膜（Lagrangian brane），$\Hom(\B_\op,\B_\op)$ 就是 $\bar\partial_\cA$ 上同调群 $H^i(L_\op, \O)$ 的直和。第 \ref{dualco} 节已经说明，$L_\op$ 非典范地同构于一个向量空间。因此当 $i>0$ 时上同调 $H^i(L_\op,\O)$ 为零，而 $\Hom(\B_\op,\B_\op)$ 就是 $L_\op$ 上全纯函数的（未形变的！）交换代数。\zy{原文此处作 $\Hom(L_\op,L_\op)$，按上下文应为 $\Hom(\B_\op,\B_\op)$；译文径改。}于是，与 $B$-模型的对偶为“$\A$ 必须交换”提供了另一种解释。此外它还表明，代数 $\A$ 的“谱”——这里取其抽象意义，即其 1 维复表示所构成的空间——是参数化 oper 的“簇”\footnote{对我们而言，“簇”不过是“复流形”的同义词。} $L_\op$，正如\cite{BD} 最初所示。用前几节的语言说，这就是微分算子 $\D_\P$ 与函数 $T_\P$ 之间的一个典范等同。完全相同的论证表明，$\b\A=\Hom(\B_\bop,\B_\bop)$ 是反全纯 oper 的簇 $L_\bop$ 上全纯函数的代数。\zy{原文此处作 $\Hom(\B_\op,\B_\op)$，按上下文（$\b\A$ 的对偶描述）应为 $\Hom(\B_\bop,\B_\bop)$；译文径改。}

如第 \ref{dualco} 节所解释，oper 的簇非典范同构于 Hitchin 纤维化的底空间。因此，“$\A$ 是 $L_\op$ 上全纯函数的代数”这一事实，可视为“$\A_0$ 是 Hitchin 纤维化底空间上全纯函数的代数”这一事实的某种量子形变。对 $\b\A$ 当然也有类似陈述。

然而，我们并不想在已指出的抽象意义上理解 $\A\times \bA$ 的谱，而是把它看成 $\H=\Hom(\bB_\cc,\B_\cc)$ 上的具体算子。对偶理论则用 $\H=\Hom(\B_\bop,\B_\op)$ 给出对偶描述。若 $G^\vee$ 的中心平凡，就可以只在 $\sigma$-模型（而不是带 $\zZ(G^\vee)$ 规范场的 $\sigma$-模型）中分析此事。我们先考虑这一情形。事情之所以简单，是因为所涉及的膜都是秩 1 的拉格朗日膜，支撑在复拉格朗日子流形 $L_\op$ 与 $L_\bop$ 上。分析这一问题时，我们假定 $L_\op$ 与 $L_\bop$ 只在孤立点处横截相交。\zy{“横截（transverse）”指两子流形在交点处的切空间之和等于整个切空间；“孤立”指交点集离散、没有连续族。二者合起来意味着 $L_\op$ 与 $L_\bop$ 在每个交点附近只相交一次，这正是把交点作为 $\H$ 的一组基来计数的前提。译者按（无外部依据）。} 这一点对 $\SL(2,\C)$ 已知成立；一般而言，它是 Etingof、Frenkel 与 Kazhdan 的猜想之一\cite{EFK,EFK2}。交点孤立且横截，实际上是这一对偶的一项预言；它是 $\H$ 上的厄米（Hermitian）形式为正定所必需的，而正定性是 $A$-模型构造所预期的——在该构造中，$\H$ 是 $\lmark^2$ 半密度的希尔伯特空间。\zy{半密度（half-density）可看作密度的平方根：两个半密度相乘得到通常的密度，因此无需预先选定测度即可积分。$\lmark^2$ 半密度空间因而带有典范的厄米内积 $(s,t)=\int \bar s\,t$，这就是几何量子化中的“典范希尔伯特空间”。参见 nLab 条目 “density” 与 “canonical Hilbert space of half-densities”。} 遗憾的是，要解释这一点，需要对 $B$-模型量子力学作相当详细的讨论，这部分内容已放入附录 \ref{bmodel}。（在该附录中我们会看到，正定性实际上还需要另一个尚未证明的必要条件。）

设 $\Upsilon=L_\op\cap L_\bop$。假定交点孤立且横截，$\H=\Hom(\B_\bop,\B_\op)$ 就有一组基：对每个 $u\in \Upsilon$ 恰有一个基向量 $\psi_u$。这是关于 $B$-模型中拉格朗日膜之交的一般性陈述。具体地说，$L_\op$ 是 $\M_H(G,C)$ 中参数化“作为全纯 oper 的平坦丛”的子簇，$L_\bop$ 是参数化“作为反全纯 oper 的平坦丛”的子簇；由此可见，一个交点代表一个平坦 $G_\C^\vee$ 丛，它同时是全纯 oper 又是反全纯 oper。

我们回顾，$\Hom(\B_\bop,\B_\op)$ 上厄米形式的定义用到一个反全纯对合 $\tau$，其作用为 $(A,\phi)\to (A,-\phi)$。因此 $\tau$ 把复平坦联络 $\cA=A+\i\phi$ 变为复共轭平坦联络 $\b\cA=A-\i\phi$。回忆 $A$ 是规范场，而该理论的规范群是紧致形式 $G^\vee$。在数学上，$G^\vee_\C$ 中固定 $G^\vee$ 的对合是复共轭与 Chevalley 对合之积，故 $\tau$ 以该乘积作用在 $\cA$ 上（至多差一个内自同构，这里指 $G^\vee$ 值的规范变换）。\zy{Chevalley 对合是约化群的一个标准对合：它在极大环面上取逆 $t\mapsto t^{-1}$，并把每个根子群送到相对的负根子群。取定 $G^\vee_\C$ 的分裂实形式后，复共轭固定该分裂形式；它与 Chevalley 对合之积才固定紧致实形式 $G^\vee$，此即原文此句的由来。译者按（无外部依据）。}

在态空间 $\H$ 的折叠构造中，$\tau$ 通过交换两个因子而反全纯地作用在 $\h Y=Y_1\times Y_2$ 上。这意味着，在展开构造中，$\tau$ 交换图 \ref{folding}(b) 中带条的两端。不难具体看出其缘由。若 $\cA$ 是作为全纯 oper 的复平坦联络，则 $\b\cA$ 是作为反全纯 oper 的复平坦联络；类似地，若 $\cA$ 是反全纯 oper，则 $\b\cA$ 是全纯 oper。因此 $\tau$ 交换 $L_\op$ 与 $L_\bop$，并同样交换\footnote{$\tau$ 交换 $\B_\op$ 与 $\B_\bop$ 这一陈述，是在底层的物理 $\sigma$-模型中成立的。由于 $\tau$ 在相关复结构下反全纯地作用在 $\M_H(G^\vee,C)$ 上，它把 $B$-模型与一个共轭 $B$-模型互换，因而并不是 $B$-模型的对称性（$\B_\op$ 与 $\B_\bop$ 在 $B$-模型及其共轭模型中都是合法的膜）。用于定义厄米结构（式 (\ref{hermdef})）并交换 $\B_\op$ 与 $\B_\bop$ 的 $B$-模型对称性是 $\Theta_\tau=\Theta\tau$，其中 $\Theta=\sf{CPT}$。} $\B_\op$ 与 $\B_\bop$。

设点 $u\in L_\op\cap L_\bop$ 对应一个复平坦丛 $E^\vee_\C$，它同时是全纯 oper 与反全纯 oper。那么其复共轭 $\b E^\vee_\C$ 也同时是全纯 oper 与反全纯 oper。若 $\b E^\vee_\C$ 作为平坦丛不与 $E^\vee_\C$ 规范等价，则 $\b E^\vee_\C$ 对应一个不同于 $u$ 的点 $\b u\in L_\op\cap L_\bop$。若如此，$u$ 与 $\b u$ 将对应 $\H$ 中不同的基向量 $\psi_u$ 与 $\psi_{\b u}$，而且它们会被 $\Theta_\tau$ 互换。$B$-模型的自然配对在交点基下是对角的：基向量可归一化，使得对 $u,u'\in \Upsilon$ 有 $(\psi_u,\psi_{u'})=\delta_{uu'}$。因此，若 $\Theta_\tau$ 互换两个不同的基向量 $\psi_u$ 与 $\psi_{\b u}$，则对式 (\ref{hermdef}) 所定义的厄米内积而言，$\psi_u$ 与 $\psi_{\b u}$ 都是迷向向量（null vector）。对偶预言这一厄米内积应当是正定的，因为在 $A$-模型一侧，$\H$ 由量子化一个余切丛得到，是半密度的希尔伯特空间。于是我们预期：一个同时为全纯与反全纯 oper 的平坦丛实际上是实的。这一点是\cite{EFK,EFK2} 中的猜想，并对 $G^\vee=\U(1)$ 通过一个显式（但出人意料地不平凡）的计算得到证明；对 $G^\vee=\SU(2)$ 结果也已知\cite{Fal,Go}。

\begin{yizhu}[译者按：正定性与实 oper]
这段论证的链条容易读反，梳理如下。$\tau$ 交换 $L_\op$ 与 $L_\bop$；若 $u$ 与 $\b u$ 是不同的交点，$\Theta_\tau$ 就把基向量 $\psi_u$ 与 $\psi_{\b u}$ 互换（至多差一个归一化因子）。由式 (\ref{hermdef})，$\langle\psi_u,\psi_u\rangle=(\Theta_\tau\psi_u,\psi_u)$。右端的两个因子对应不同的交点，由 $B$-模型配对在交点基下的对角性，其值为零；故 $\psi_u$ 与 $\psi_{\b u}$ 都是迷向向量。正定性不允许非零迷向向量，于是只能有 $u=\b u$，即该丛是实的。须注意：原文把正定性说成对偶的一项预言（在 $A$-模型一侧由构造预期），附录 \ref{bmodel} 还指出存在另一个未证明的必要条件，故这里是条件性的论证，不宜当作已证定理。（本注仅为对原文论证的梳理，无外部依据。）
\end{yizhu}

最后，我们可以利用对偶来预言全纯与反全纯 Hitchin 哈密顿量作为 $\H$ 上算子的谱。设 $H_{\P,\alpha}$ 是一个量子 Hitchin 哈密顿量，即 $\A=\Hom(\B_\cc,\B_\cc)$ 中的一个元素。对偶把 $\A$ 与 $\Hom(\B_\op,\B_\op)$ 等同，从而把 $H_{\P,\alpha}$ 与 $L_\op$ 上的全纯函数 $f_{\P,\alpha}$ 等同。作用在对应点 $u\in L_\op\cap L_\bop$ 的基向量 $\psi_u$ 上时，$H_{\P,\alpha}$ 就是以相应值 $f_{\P,\alpha}(u)$ 作乘法。类似地，若 $H_{\b\P,\b\alpha}\in \b\A=\Hom(\b\B_\cc,\b\B_\cc)$ 是一个反全纯的量子 Hitchin 哈密顿量，则它在对偶下对应 $L_\bop$ 上的全纯函数 $f_{\bar\P,\b\alpha}$，并在 $\psi_u$ 上以 $f_{\b\P,\b\alpha}(u)$ 作乘法。至此完成了量子 Hitchin 哈密顿量本征值的描述。

\subsection{含中心的情形}\label{inccenter}

要修改这一描述以计入 $G_\C^\vee$ 的中心，并不困难。照例考虑 $M=\Sigma\times C$ 上的 $B$-模型。它局域化到 $\Sigma\times C$ 上的平坦丛。在我们的应用中，$\Sigma=\R\times I$ 可缩，因此 $M$ 上的平坦丛是 $C$ 上某个平坦丛的拉回。当带条两端取 oper 与反 oper 边界条件时，$C$ 上的平坦丛同时是全纯 oper 与反全纯 oper；因此它是一个实 oper。\zy{“实 oper”指该平坦丛在反全纯对合 $\tau$ 下保持不变，即它与自己的复共轭 $\b E^\vee_\C$ 规范等价；等价地，$L_\op\cap L_\bop$ 中的对应点满足 $\b u=u$。用正文的话说，它同时是全纯 oper 与反全纯 oper。注意“实”指该平坦丛具有实结构，并不意味着它来自紧致形式 $G^\vee$ 的平坦丛。译者按（无外部依据）。}

oper 丛无论是否实，都是不可约的；其自同构群只由规范群的中心 $\zZ(G^\vee)$ 构成。然而，如第 \ref{opergauge} 节所解释，这些边界条件还在两个边界上给出 $\zZ(G^\vee)$ 规范对称性的平凡化 $s_\ell$ 与 $s_r$，但须模去按 $(s_\ell, s_r)\to (b s_\ell,b s_r)$（$b\in \zZ(G^\vee)$）作用的规范变换。对给定对应点 $u\in \Upsilon$ 的一个实 oper，设 $\T_u$ 为配对 $s_\ell,s_r$ 在 $\zZ(G^\vee)$ 作用下的等价类集合。$B$-模型局域化到由点 $u,\varepsilon$ 构成的孤立集合上，其中 $u\in \Upsilon$，$\varepsilon\in \T_u$。因此，中心非平凡的一般情形下，希尔伯特空间 $\H$ 有一组基 $\psi_{u,\varepsilon}$，指标取遍上述 $u,\varepsilon$。\zy{$\T_u$ 是一个 $\zZ(G^\vee)$-torsor（主齐性空间）：$\zZ(G^\vee)$ 在其上自由且传递地作用，但没有典范的单位元。中心平凡时 $\T_u$ 退化为单点，基 $\psi_{u,\varepsilon}$ 就回到此前单个交点对应的 $\psi_u$；中心非平凡时同一 $u$ 会带出多个 $\varepsilon$，它们给出相同的 Hitchin 本征值而增大重数。参见 nLab 条目 “torsor”。}

可以把 $\varepsilon\in\T_u$ 看成带条左右边界之间的一种整体和乐（global holonomy）。这一借助 torsor $\T_u$ 的精细化，在通过 $L_\op$ 或 $L_\bop$ 上的全纯函数对代数 $\A$ 与 $\bA$ 作对偶描述时并不重要，因为每个代数只作用在带条的一侧。在确定 Hitchin 哈密顿量的本征值时，它同样不重要：后者只依赖于把 $\A$ 与 $\bA$ 解释为 $L_\op$ 与 $L_\bop$ 上的函数，而不敏感于横穿带条的整体和乐。它确实影响本征值的重数，因为 $u$ 相同而 $\varepsilon$ 不同的本征向量 $\psi_{u,\varepsilon}$ 具有相同的 Hitchin 哈密顿量本征值。并且在描述 't Hooft 或 Hecke 算符的本征值时，它也会起作用——我们接下来就转向这一点。
